% =====================================================================
\kicker{Guide topic 4}
\section{Application Definition}\label{sec:app}
% =====================================================================
\resultstrip{Define the input, processing and output of the rack controller and its emulator interface.}
  {8 commands, 7 decision rules, 12 emulator runs}
  {Tables~\ref{tab:commands}–\ref{tab:emuruns}}

\subsection{System context}\label{sec:context}
The controller protects one server rack in a Johor data hall, one of the fastest-growing data-centre locations in the region. The temperature thresholds follow the ASHRAE thermal guideline for data-processing environments~\cite{ashrae}: the recommended inlet range ends at \textbf{27\,°C} (Warm from this value) and the class A1 allowable maximum is \textbf{32\,°C} (Critical above it). The rack is fed by a 32\,A, 230\,V circuit. With the usual limit of 80\,\% of a breaker's rating for a continuous load, the warning level is $I_{\text{warn}} = 0.8 \times 32\,\mathrm{A} = \mathbf{25.6\,A}$, and a current above the 32\,A rating trips the EPO.

\subsection{System input}\label{sec:input}
Readings arrive as one 6-byte \textbf{SAMPLE} packet: header \hex{0x5A}, voltage in 1\,V steps, current in 0.2\,A steps and temperature in 0.5\,°C steps, then the CRC (example below). Configuration and status use 4-byte commands, \code{[CMD]\allowbreak[DATA]\allowbreak[CRC\_LO]\allowbreak[CRC\_HI]} (\tabref{tab:commands}). Both formats end with a CRC-16-CCITT (polynomial \hex{0x1021}, seed \hex{0xFFFF}, the CCSDS frame check~\cite{ccsds}) over the preceding bytes. P0 is the operator reset button; P1–P3 are not used.

\smallskip
\noindent\begin{tikzpicture}[x=1mm,y=1mm]
  \foreach \i/\n/\v/\m/\c in {0/HEAD/5A/SAMPLE/navy,1/V/E6/230\,V/navy,2/I/7F/25.4\,A/navy,3/T/41/32.5\,°C/navy,4/CRC\_LO/43/{}/river,5/CRC\_HI/CB/{}/river}{
     \draw[\c, line width=0.7pt, fill=\ifnum\i<4 white\else riversoft\fi] ({\i*15.9},3) rectangle ({\i*15.9+14.6},14);
     \node[font=\sffamily\tiny\bfseries, text=\c] at ({\i*15.9+7.3},11.8) {\n};
     \node[font=\ttfamily\small, text=ink] at ({\i*15.9+7.3},8.2) {0x\v};
     \node[font=\sffamily\tiny, text=slate] at ({\i*15.9+7.3},4.7) {\m};}
  \draw[->, slate, line width=0.5pt] (0,0.8) -- (94,0.8);
  \node[lbl, anchor=west, fill=white] at (38,0.8) {sent first\enspace→\enspace sent last};
  \node[kicker, anchor=west] at (100,13) {PACKET ANATOMY:\enspace CASE C03, 230\,V · 25.4\,A · 32.5\,°C};
  \node[lbl, anchor=west] at (100,9.3) {CRC-16 of the first four bytes\,=\,\hex{0xCB43}, low byte first};
  \node[font=\ttfamily\small\bfseries, text=river, anchor=west] at (100,4.8) {-t4 0x417FE65A -t2 0xCB43};
  \node[lbl, anchor=west] at (100,1.2) {the same packet as \texttt{./emu} options (\secref{sec:emu})};
\end{tikzpicture}

% Table 2 -- Command set of dc_teammate.s (.equ block and dispatch chain)
\begin{table}[H]
\caption[Command set]{Command set. Replies: \texttt{0xA5} accepted, \texttt{0x5A} bad CRC (nothing executed), \texttt{0x3C} invalid or unknown.}
\label{tab:commands}
{\tablestyle
\begin{tabularx}{\linewidth}{@{}P{15mm}P{35mm}P{37mm}P{14mm}L@{}}
\toprule
\colhead{CMD} & \colhead{Name} & \colhead{Data} & \colhead{Reply} & \colhead{Effect}\\
\midrule
\hex{0x5A} & SAMPLE & V 180–254; I 0–254; T 1–254 & A5 / 3C & Store V and I; T\,×\,5 (\code{mul}) → 0.1\,°C; classify; current warning; EPO trip\\
\hex{0x40} & STAGE\_HIGH & high byte & A5 & Stage the high byte of a temperature threshold\\
\hex{0x41} / \hex{0x43} & SET\_WARM / SET\_CRIT & low byte & A5 / 3C & Commit (0.1\,°C units) only after STAGE\_HIGH and only if Warm < Critical\\
\hex{0x46} / \hex{0x48} & SET\_WARN\_CURRENT / SET\_TRIP\_CURRENT & raw 1 to trip$-$1 / above warning, not 255 & A5 / 3C & Current warning level (default 128 = 25.6\,A) / trip level (default 160 = 32.0\,A)\\
\hex{0x50} & STATUS & 0 & 4 bytes & Flags (class, warning, EPO, temperature received, staged), temperature, current, EPO trips\\
\hex{0x52} & COUNTERS & 0 & 4 bytes & CRC rejects, then three zero bytes\\
other & any other value & — & 3C & Unknown command: nothing executed\\
\bottomrule
\end{tabularx}}
\end{table}


\subsection{System processing and output}\label{sec:processing}
\tabref{tab:decisions} is the complete rule set. A packet with a bad CRC or out-of-range data receives its reply (\hex{0x5A} or \hex{0x3C}) and changes no state. Classification uses the latest SAMPLE; a reading of exactly 32.0\,°C remains Warm, and exactly 32.0\,A remains a warning. Each complete packet receives one reply.

% Table 3 -- Decision table: every rule of Sections 4.1 and 4.3 in one place
% (cmd_sample_process, check_reset, update_outputs in dc_teammate.s).
\newcommand{\swatch}[1]{\tikz[baseline=-0.6ex]\fill[#1] (0,-0.9ex) rectangle (0.9ex,0.6ex);\hspace{0.35em}}
\begin{table}[H]
\caption[Decision table: conditions, triggers and pins]{Decision table: each condition the firmware evaluates, its trigger and basis, the pins it drives and the value \texttt{./emu -o} reads. Pin values add, e.g.\ \texttt{0xE0} = fans max + warning + EPO.}
\label{tab:decisions}
{\tablestyle
\begin{tabularx}{\linewidth}{@{}P{30mm}LP{50mm}P{11mm}@{}}
\toprule
\colhead{Condition} & \colhead{Trigger and basis} & \colhead{Pins driven} & \colhead{\texttt{-o}}\\
\midrule
\swatch{teal}Normal & inlet below 27.0\,°C, inside the ASHRAE recommended range & none & \hex{0x00}\\
\swatch{amber!45}Warm & 27.0\,°C ≤ inlet ≤ 32.0\,°C & P4 fan boost & \hex{0x10}\\
\swatch{alarm}Critical temperature & inlet above 32.0\,°C, the class A1 allowable maximum & P5 fans at maximum + P7 EPO & \hex{0xA0}\\
\swatch{amber}Current warning & I ≥ 25.6\,A (80\,\% of the 32\,A feed); follows the latest reading & P6 current warning & \hex{0x40}\\
\swatch{alarm}Current trip & I above 32.0\,A, the feed rating & P6 warning + P7 EPO & \hex{0xC0}\\
\swatch{alarm}EPO latched & stays on after its cause is gone, until a valid reset & P7 EPO (+ current P4–P6) & \hex{0x80}\\
\swatch{teal}EPO reset & new P0 press while the inlet is Normal and I < 25.6\,A & releases P7; a held button never does & —\\
\bottomrule
\end{tabularx}}
\end{table}


\subsection{Emulator values (./emu)}\label{sec:emu}
\noindent\begin{minipage}[t]{0.4\linewidth}
\vspace{0pt}\RaggedRight
The organizer's emulator~\cite{repos} drives the pins and the UART of the FPGA top. Each run resets the processor and stores at most 8 reply bytes. Every register access waits 50\,ms, so consecutive UART bytes reach the chip at least 100\,ms apart. \tabref{tab:emuruns} lists the 12 demonstration runs, which use the options on the right.
\end{minipage}\hfill
\begin{minipage}[t]{0.575\linewidth}
\vspace{0pt}
{\tablestyle\rowcolors{1}{white}{mist!70}
\begin{tabularx}{\linewidth}{@{}P{17mm}L@{}}
\toprule
\emu{-d 0xF0} & P7–P4 observed by the emulator, P3–P0 driven by it\\
\emu{-i V} & Drive P3–P0: \hex{0x01} presses P0, \hex{0x00} releases it\\
\emu{-o V} & Compare P7–P4 with V\\
\emu{-t4}, \emu{-t2} & Send 4 or 2 bytes, least-significant first: SAMPLE \emu{-t4 0x\{T\}\{I\}\{V\}5A -t2 0x\{CRC\}}; command \emu{-t4 0x\{CRC\}\{DATA\}\{CMD\}}\\
\emu{-rN V} & Expect N reply bytes; the first byte is the least significant\\
\bottomrule
\end{tabularx}}
\end{minipage}

% Table 5 -- The 12 verified ./emu runs (command lines verbatim, as printed by
% tb_dc_teammate.v; identical to teammate_version/emu_runs.txt).
\newcommand{\emuline}[1]{{\ttfamily\fontsize{7}{8.2}\selectfont\RaggedRight
  \exhyphenpenalty=10000 \hyphenpenalty=10000
  \automatichyphenpenalty=10000 \explicithyphenpenalty=10000 #1\par}}
\begin{table}[!t]
\caption[Verified ./emu demonstration runs]{Emulator command lines, printed by the testbench after each case passed; \figref{fig:matrix} decodes each step. Example: C03 STATUS \texttt{0x017F414B} = flags \texttt{0x4B} (Critical, EPO, temperature received), 32.5\,°C (\texttt{0x41}), 25.4\,A (\texttt{0x7F}), 1 EPO trip.}
\label{tab:emuruns}
{\tablestyle\renewcommand{\arraystretch}{1.08}
\begin{tabularx}{\linewidth}{@{}P{9mm}L@{}}
\toprule
\colhead{Case} & \colhead{\texttt{./emu} command line (verified)}\\
\midrule
C01 & \emuline{./emu -d 0xF0 -t4 0x357FE65A -t2 0xF550 -r1 0xA5 -o 0x00 -t4 0x13B00050 -r4 0x007F3540}\\
C02 & \emuline{./emu -d 0xF0 -t4 0x367FE65A -t2 0xC533 -r1 0xA5 -o 0x10 -t4 0x407FE65A -t2 0xDB62 -r1 0xA5 -o 0x10 -t4 0x13B00050 -r4 0x007F4041}\\
C03 & \emuline{./emu -d 0xF0 -t4 0x417FE65A -t2 0xCB43 -r1 0xA5 -o 0xA0 -t4 0x13B00050 -r4 0x017F414B}\\
C04 & \emuline{./emu -d 0xF0 -t4 0x327FE65A -t2 0x85B7 -r1 0xA5 -o 0x00 -t4 0x3280E65A -t2 0x8648 -r1 0xA5 -o 0x40 -t4 0x32A0E65A -t2 0x80AE -r1 0xA5 -o 0x40}\\
C05 & \emuline{./emu -d 0xF0 -t4 0x32A1E65A -t2 0xB39F -r1 0xA5 -o 0xC0 -t4 0x13B00050 -r4 0x01A1324C}\\
C06 & \emuline{./emu -d 0xF0 -t4 0x41A1E65A -t2 0xFD6B -r1 0xA5 -o 0xE0 -t4 0x32A1E65A -t2 0xB39F -r1 0xA5 -o 0xC0 -t4 0x417FE65A -t2 0xCB43 -r1 0xA5 -o 0xA0}\\
C07 & \emuline{./emu -d 0xF0 -t4 0x41A1E65A -t2 0xFD6B -r1 0xA5 -o 0xE0 -t4 0x3680E65A -t2 0xC6CC -r1 0xA5 -o 0xD0 -t4 0x357FE65A -t2 0xF550 -r1 0xA5 -o 0x80 -i 0x01 -o 0x00 -i 0x00}\\
C08 & \emuline{./emu -d 0xF0 -t4 0x417FE65A -t2 0xCB43 -r1 0xA5 -o 0xA0 -t4 0x407FE65A -t2 0xDB62 -r1 0xA5 -o 0x90 -t4 0x357FE65A -t2 0xF550 -r1 0xA5 -o 0x80 -i 0x01 -o 0x00 -i 0x00}\\
C09 & \emuline{./emu -d 0xF0 -t4 0x32A1E65A -t2 0xB39F -r1 0xA5 -o 0xC0 -t4 0x3280E65A -t2 0x8648 -r1 0xA5 -o 0xC0 -i 0x01 -o 0xC0 -i 0x00 -t4 0x327FE65A -t2 0x85B7 -r1 0xA5 -o 0x80 -i 0x01 -o 0x00 -i 0x00}\\
C10 & \emuline{./emu -d 0xF0 -t4 0x417FE65A -t2 0xCB43 -r1 0xA5 -o 0xA0 -t4 0x007FE65A -t2 0x93A6 -r1 0x3C -o 0xA0 -t4 0x357FE65A -t2 0xF551 -r1 0x5A -o 0xA0}\\
C11 & \emuline{./emu -d 0xF0 -t4 0x417FE65A -t2 0xCB43 -r1 0xA5 -o 0xA0 -i 0x01 -o 0xA0 -i 0x00 -t4 0x367FE65A -t2 0xC533 -r1 0xA5 -o 0x90 -i 0x01 -o 0x90 -i 0x00 -t4 0x357FE65A -t2 0xF550 -r1 0xA5 -o 0x80 -i 0x01 -o 0x00 -i 0x00}\\
C12 & \emuline{./emu -d 0xF0 -i 0x01 -t4 0x32A1E65A -t2 0xB39F -r1 0xA5 -o 0xC0 -t4 0x327FE65A -t2 0x85B7 -r1 0xA5 -o 0x80 -i 0x00 -i 0x01 -o 0x00 -i 0x00}\\
\bottomrule
\end{tabularx}}
\end{table}

